\section{State of the art}

In this section realizability and implementability will be used as the same problem.
Informally, an MSC can be defined realizable if a distributed 
implementation exists that exactly reproduces the graph's behaviors.

\subsection{Realizability by Alur}

\subsubsection{Definitions}

\begin{definition}[Weakly-implies]
A set $L$ of MSCs weakly implies an MSC $M$ iff for all $1\leq i \leq n$, 
there exists an MSC $M_i \in L$ such that $M|_i = M_i|_i$.
\end{definition}

\begin{definition}{Weakly realizable}
The weak closure $L^w$ of a set $L$ of MSCs contains all the MSCs $L$ weakly implies.
The set $L$ is weakly realizable iff $L=L^w$. The MSC-graph $G$ is said to be weakly
realizable if the set $L(G)$ of MSCs is. Thus, a weakly realizable graph already contains
all the implied scenarios.
\end{definition}

\begin{definition}[Deadlock-free]
Consider a set $A_i$ of concurrent automata and the product $A =\prod_i A_i$
A state $q$ of the product $A$ is said to be a deadlock state if no accepting state of $A$ is reachable from $q$.
The product $A$ is said to be deadlock-free if no state reachable from its initial state is a deadlock state.
\end{definition}

\begin{definition}[Safely realizable]
A set $L$ of MSCs is said to be safely realizable if $L=L\prod A_i$ for some $\langle A_i | 1\leq i \leq n \rangle$
such that $\prod A_i$ is deadlock-free.
\end{definition}

\subsubsection{Realizability problem}

In \cite{alur2005realizability}, they prove many interesting results within the scope of MSC graphs.
They first examine the realizability of bounded MSC-graphs. Two types are considered: 
\textit{weak} (allowing deadlocks) and \textit{safe} (deadlock-free). While weak realizability is \verb|coNP|-complete
and safe realizability solvable in polynomial time for finite MSC sets, they show that for
bounded MSC-graphs, weak realizability becomes undecidable, and safe realizability lies in 
\verb|EXPSPACE|.

\begin{theorem}
Given a bounded MSC graph G, checking if G is weakly realizable is undecidable.
\end{theorem}

\begin{proof}
The proof is a reduction from the post correspondence problem (PCP). In the paper \cite{alur2005realizability}
they further reduce the problem to a relaxed version. First, they reduce PCP to OneSidedPCP and then to
Relaxed PCP (RPCP). We skip this part as is not meaningful for the proof. We can state the RPCP to be the following
problem: given $\{(v_1, w_1), (v_2, w_2), ..., (v_r, w_r)\}$, determinate whether there are indices $i_1, ..., i_m$
such that $x_{i_1}...x_{i_n}=y_{i_1}...y_{i_m}$, where $x_{i_j}, y_{i_j} \in \{v_{i_j},w_{i_j}\}$, for some
index $i_l \  x_{i_l}\neq y_{i_l}$, and for all $j\leq m,y_{i_1}...y_{i_j}$ is a prefix of $x_{i_1}...x_{i_j}$.
The relaxed problem essential says that must have the left side of the equivalence that grows more than the right side,
we must have also some tuple that are different, and that we can choose whatever side of tuple for the composition
of the string. Now we reduce RPCP to weak realizability. 

\newcommand{\syncmess}[3]{%
  \mess{#1}{#2}{#3}%
  \mess{}{#3}{#2}%
  \nextlevel
}

\newcommand{\asyncmess}[3]{%
  \mess{#1}{#2}{#3}%
  \nextlevel
}

\begin{figure}[h]
\centering
\begin{msc}[draw frame=none, draw head=none, msc keyword=, head height=0px, label distance=0.5ex, foot height=0px, foot distance=0px]{}
    \declinst{P1}{P1}{}
    \declinst{P2}{P2}{}
    \declinst{P3}{P3}{}
    \declinst{P4}{P4}{}

    \syncmess{$(i,0)$}{P1}{P2}
    \asyncmess{$i$}{P1}{P4}
    \syncmess{$(i,0)$}{P3}{P4}
    \syncmess{$v_i^1$}{P2}{P3}
    \syncmess{...}{P2}{P3}
    \syncmess{$v_i^c$}{P2}{P3}
\end{msc}
\caption{The $M_i^0$ MSC. For a string $u$, let $u_l$ denote the $l$’th character of the string.}
\end{figure}

\begin{figure}[h]
\centering
\begin{msc}[draw frame=none, draw head=none, msc keyword=, head height=0px, label distance=0.5ex, foot height=0px, foot distance=0px]{}
    \declinst{P1}{P1}{}
    \declinst{P2}{P2}{}
    \declinst{P3}{P3}{}
    \declinst{P4}{P4}{}

    \syncmess{$(i,1)$}{P1}{P2}
    \asyncmess{$i$}{P1}{P4}
    \syncmess{$(i,1)$}{P3}{P4}
    \syncmess{$w_i^1$}{P2}{P3}
    \syncmess{...}{P2}{P3}
    \syncmess{$w_i^d$}{P2}{P3}
\end{msc}
\caption{The $M_i^1$ MSC. For a string $u$, let $u_l$ denote the $l$’th character of the string.}
\end{figure}


\textit{Encodig:} Given a finite set $L$ of
MSCs, let $L^*$ denote the MSC graph that consists of the complete
graph with $|L|$ vertices, one per MSC in the set $L$, dummy initial
and terminal vertices $v_I$, $v_T$ with empty MSCs, and edges from
$v_I$ to all vertices of $L$ and from those to $v_T$. Thus, an MSC of
this graph is simply a concatenation of MSCs from the set $L$.

In the sequel, we say that a process $p$ synchronously sends a
message $m$ to process $q$, if $p$ sends $m$ to $q$ immediately
followed by $q$ sending the message $m$ back to $p$. In figures,
such messages will be depicted by double arrows.

Given an instance $\Delta = \{(v_1, w_1), \ldots, (v_m, w_m)\}$ of
RPCP, we build a set $L$ of MSCs over 4 processes as follows.
For each pair $(v_i, w_i)$ we build two MSCs $M^0_i$ and $M^1_i$,
which are depicted in Fig.~3.
Thus in $M^0_i$, process 1 sends synchronously $(i, 0)$ to process 2,
then sends the index $i$ to process 4, and then process 4 sends
synchronously $(i, 0)$ to process 3. After that, process 2
synchronously sends the sequence of characters of $v_i$ to process 3 (note we assume $c$
is the length of $v_i$ and $d$ the length of $w_i$ in the figure), $M^1_i$
is similar. Observe that the communication graph of each of these
MSCs is strongly connected and involves all the processes, and
hence, the MSC graph $L^*$ is bounded.

\textbf{Claim 1:} $\Delta \in \text{RPCP}$ iff $L^*$ is not weakly realizable.

\textit{Proof.} For the ``only if" direction, suppose $R = (i_1, a_1, b_1, i_2, a_2, b_2, \ldots,
i_m, a_m, b_m)$ are the indices for a solution to $\mathcal{B}$, and the bits $a_j$ and $b_j$
indicate which string ($v_{i_j}$ or $w_{i_j}$) is chosen to go into the two (left and right)
long strings.

Consider the new MSCs $M$ and $M'$ obtained from the sequences
$M = M^{a_1}_{i_1} \cdots M^{a_m}_{i_m}$ and $M' = M^{b_1}_{i_1} \cdots M^{b_m}_{i_m}$.
Executions of both of these (sequences of) MSCs must exist in any
realization of $L^*$. We then look at the projections $M|_1$, $M|_2$, $M|_3$,
and $M|_4$ of $M$, and $M'|_1$, $M'|_2$, $M'|_3$, $M'|_4$ of $M'$ onto the
4 processes. Now consider an MSC $M''$ formed from $M'|_1$, $M'|_2$,
$M|_3$, and $M|_4$. The claim is that the combined MSC $M''$ is weakly
implied by $L^*$. By definition, the only thing to establish is that $M''$
is indeed an MSC, in the sense that it is acyclic, well-formed, and complete.

The only new situation in terms of communication in $M''$ is the
communication between $P_1$ and $P_4$, and between $P_2$ and $P_3$.
But the communication between $P_1$ and $P_4$ is consistent in
$M'|_1$ and $M|_4$ (i.e., the sequence of messages sent from $P_1$ to
$P_4$ in $M'|_1$ is equal to the sequence of messages received in $M|_4$),
and the communication between $P_2$ and $P_3$ is consistent in
$M'|_2$ and $M|_3$ because $R$ is a solution to the RPCP.

Furthermore, the acyclicity of $M''$ follows from the property of the
solution that the string formed by the first $j$ words on processes 1
and 2 is always a prefix of the string formed by the first $j$ words
on processes 3 and 4. Consequently, each message from $P_1$ to $P_4$
is sent before it needs to be received.

But note that $M''$ cannot itself be in $L^*$ because there must be
some index $i_j$ where $a_j \neq b_j$, and no MSC exists in $L$ where,
after process 1 announces the index, what process 2 sends is not
identical to what process 3 receives.

Now, for the “if” direction, suppose there is some MSC $M^@$ which
exists in any realization of $L^*$, but is not in $L^*$ itself. We want
to derive a solution to $\mathcal{B}$ from $M^@$.

First, it is clear that the projection $M^@|_1$ must consist of a sequence
of pairs of messages (the first of each pair acknowledged), sent from
process 1 to processes 2 and 4, respectively, with messages $(i, b)$ and $i$.

Likewise, it is clear that, in order for process 2 to receive those messages,
$M^@|_2$ must consist of a sequence of receipts of $(i, b)$ pairs, and after
each $(i, b)$, either $v_i$ or $w_i$ is sent to process 3, based on whether
$b = 0$ or $b = 1$, before the next index pair is received.

Likewise, $M^@|_4$ consists of a sequence of receipts of index $i$ from
process 1, followed by sending of $(i, 0)$ or $(i, 1)$ to process 3, and
$M^@|_3$ consists of a sequence of receipt of $(i, 0)$ or $(i, 1)$ followed
by receipt of $v_i$ or $w_i$, respectively.

Now, since $M^@$ is not in $L^*$, for some index $i$ the choice of $v_i$ or
$w_i$ must differ on process 2 and process 3. (Note, we are assuming that
the buffers between processes are FIFO.)

Furthermore, because of the precedences, the prefix formed by the first
$j$ words on process 2 must precede the $(j + 1)$-th message from
process 1 to process 4, which in turn precedes the $(j + 1)$-th message
from 4 to 3, and hence the $(j + 1)$-th word on process 3. That is, the
string formed by the first $j$ words on process 2 is a prefix of the string
formed by the first $j$ words on process 3. Therefore, we can readily
build a solution for $\mathcal{B}$ from $M^@$.

\end{proof}

\subsection{Implementability by Stutz}

Part taken from \cite{stutz2024implementability}.
To start with, let $P$ be a finite set of participants and $V$ be a finite set of messages
and $\Gamma$ the set of all send and receive events. $\Gamma_p$ is the set of all send and receive events
regarding a participant $p\in P$.

\subsubsection{Definitions}

\newcommand{\send}[3]{%
  #1\rhd #2 ! #3
}

\newcommand{\recv}[3]{%
  #1\lhd #2 ? #3
}

\newcommand{\Epsilon}[0]{%
  \mathcal{E}
}

\newcommand{\doublebrace}[1]{%
  \{\!\!\{#1\}\!\!\}
}

\begin{definition}{CSM}
A \textit{communicating state machine} (CSM) is $\mathsf{A} = \doublebrace{A_p}_{p\in P}$
over $P$ and $V$ if $A_p$ is a finite state machine over $\Gamma_p$ for every $p \in P$,
denoted by $(Q_p,\Gamma_p,\delta_p,q_{0,p}, F_p)$. Let $\prod_{p\in P} Q_p$ denote
the set of global states and \verb|Chan| $=\{(p,q)| p,q \in P,p\neq q\}$ denote the set
of channels. A configuration of $A$ is a pair $(\overrightarrow{q}, \Epsilon)$, where $\overrightarrow{q}$ is a vector
of sates, one for each participant, and $\Epsilon:\text{Chan}\to V^*$ is a mapping from 
each channel to a sequence of messages. We may write $\Epsilon(p, q)$ for $\Epsilon((p, q))$.
We use $\overrightarrow{q}_p$ to denote the state of $p$ in $\overrightarrow{q}$. The
relation, denoted $\to$, is defined as follows:
\begin{itemize}
  \item $(\overrightarrow{q}, \Epsilon) \xrightarrow{\send{p}{q}{m}} (\overrightarrow{q}', \Epsilon')$
  if $(\overrightarrow{q_p},\send{p}{q}{m},\overrightarrow{q_p}') \in \delta_p, \overrightarrow{q_r} =\overrightarrow{q_r}'$
  for every participant $r\neq p$, $\Epsilon'((p,q)) = \Epsilon((p,q))\cdot m$ and $\Epsilon'(c) = \Epsilon(c)$
  for every other channel $c \in $ Chan.
  \item $(\overrightarrow{q}, \Epsilon) \xrightarrow{\recv{q}{p}{m}} (\overrightarrow{q}', \Epsilon')$
  if $(\overrightarrow{q_p},\recv{q}{p}{m},\overrightarrow{q_p}') \in \delta_p, \overrightarrow{q_r} =\overrightarrow{q_r}'$
  for every participant $r\neq p$, $\Epsilon((p,q)) = m\cdot\Epsilon'((p,q))$ and $\Epsilon'(c) = \Epsilon(c)$
  for every other channel $c \in $ Chan.
  \item $(\overrightarrow{q}, \Epsilon) \xrightarrow{\epsilon} (\overrightarrow{q}', \Epsilon)$ if 
  $(\overrightarrow{q_p},\epsilon,\overrightarrow{q_p}') \in \delta_p$ for some participant $p$, and
  $\overrightarrow{q_q} =\overrightarrow{q_q}'$ for every participant $q\neq p$.
\end{itemize}
\end{definition}

Informally, a configuration $(\overrightarrow{q}, \Epsilon)$ is a deadlock if it is not final and has no outgoing 
transitions while it is reachable if it occurs on some run of A. A CSM is deadlock-free if no reachable
configuration is a deadlock. $\Gamma^\omega$ denotes the infinite set of events.

\begin{definition}[Implementability problem]
A language $L \subseteq \Gamma^\omega$ is said to be implementable if there exists a CSM $\doublebrace{A_p}_{p\in P}$ such that
\begin{itemize}
  \item deadlock freedom: $\doublebrace{A_p}_{p\in P}$ is deadlock-free, and
  \item protocol fidelity: $L$ is the language of $\doublebrace{A_p}_{p\in P}$.
\end{itemize}
\end{definition}

